\documentclass[conference]{IEEEtran}
\IEEEoverridecommandlockouts

\usepackage[utf8]{inputenc}
\usepackage[T1]{fontenc}
\usepackage[portuguese]{babel}
\usepackage{icomma}
\usepackage{cite}
\usepackage{amsmath,amssymb,amsfonts}
\usepackage{graphicx}
\usepackage{textcomp}
\usepackage{xcolor}
\usepackage{array}
\usepackage{booktabs}
\usepackage{multirow}
\usepackage{tikz}
\usetikzlibrary{positioning,arrows.meta}
\usepackage{pgfplots}
\pgfplotsset{compat=1.16}
\usepackage{url}

\renewcommand\IEEEkeywordsname{Palavras-chave}

% Campos que o grupo tem de preencher (a vermelho). Procurar "\pendente".
\newcommand{\pendente}[1]{\textcolor{red}{\textbf{#1}}}

\definecolor{tipoum}{RGB}{31,94,153}
\definecolor{tipocore}{RGB}{120,160,205}
\definecolor{tipodois}{RGB}{214,110,30}

\begin{document}

\title{Hyper-V vs.\ VMware: Overhead da Virtualização\\
\thanks{Trabalho realizado no âmbito da Unidade Curricular de Virtualização e Centros de Dados, Departamento de Informática, Universidade da Beira Interior.}
}

\author{
\IEEEauthorblockN{Henrique Laia (M16491)}
\IEEEauthorblockA{\textit{Departamento de Informática}\\
\textit{Universidade da Beira Interior}\\
Covilhã, Portugal\\
henrique.laia@ubi.pt}
\and
\IEEEauthorblockN{Alexandre Santos (M16740)}
\IEEEauthorblockA{\textit{Departamento de Informática}\\
\textit{Universidade da Beira Interior}\\
Covilhã, Portugal\\
alexandre.minhoto.santos@ubi.pt}
\and
\IEEEauthorblockN{Daniel Basílio (E11929)}
\IEEEauthorblockA{\textit{Departamento de Informática}\\
\textit{Universidade da Beira Interior}\\
Covilhã, Portugal\\
daniel.basilio@ubi.pt}
}

\maketitle

% =================================================================== RESUMO
\begin{abstract}
Este trabalho compara o desempenho de um hipervisor nativo (Hyper-V) e de um hipervisor baseado em anfitrião (VMware Workstation~17.6) no mesmo portátil, com máquinas virtuais Ubuntu Server~26.04 iguais. Cada teste foi repetido 10 vezes e comparado com o sistema nativo. No processador, os dois hipervisores perdem 2--11\% face ao nativo. No acesso aleatório à memória, perdem 27\% (Tipo~1) e 47\% (Tipo~2). No disco, o Tipo~2 é mais rápido com um pedido de cada vez, por causa da \textit{cache} do anfitrião, mas perde 88\% de IOPS com 32 pedidos em fila. Na rede, perde 33--66\% de débito.
\end{abstract}

\begin{IEEEkeywords}
virtualização, Hyper-V, VMware Workstation, hipervisor Tipo~1, hipervisor Tipo~2, overhead, E/S
\end{IEEEkeywords}

% =============================================================== INTRODUÇÃO
\section{Introdução}
A virtualização é a base dos centros de dados e da computação em nuvem. Permite juntar vários sistemas isolados num só servidor físico, o que aproveita melhor o hardware e facilita a gestão~\cite{b3,b4}. Quem torna isto possível é o hipervisor: fica entre o hardware e os sistemas convidados e controla o acesso ao processador, à memória e aos dispositivos de Entrada/Saída (E/S). Por isso, a forma como o hipervisor está construído influencia o desempenho das máquinas virtuais.

Os hipervisores dividem-se em dois tipos~\cite{b7,b4}. O Tipo~1 (nativo ou \textit{bare-metal}) corre diretamente sobre o hardware. O Tipo~2 (\textit{host-based}) corre como uma aplicação dentro de um sistema operativo anfitrião. Costuma assumir-se que o Tipo~2 é sempre mais lento. No entanto, com a virtualização assistida por hardware, o código das máquinas virtuais já corre quase à velocidade nativa~\cite{b9,b10}, e o Tipo~2 pode até tirar partido de recursos do anfitrião, como a \textit{cache} de ficheiros. Interessa, por isso, perceber em que subsistemas há \textit{overhead} e em que situações aparece.

Neste trabalho, o Microsoft Hyper-V~\cite{b1} e o VMware Workstation~\cite{b2} foram instalados e configurados, em \textit{single node}, no mesmo portátil, e mediu-se o desempenho do processador, da memória, do disco e da rede. Os principais contributos são: (i)~uma comparação com recursos e vCPU equivalentes, em que se confirmou o modo de execução de cada hipervisor e se testou a significância das diferenças; (ii)~a explicação, com a lei de Little~\cite{b14}, do comportamento do disco, que favorece um ou outro hipervisor consoante o número de pedidos em fila; (iii)~a medição do \textit{overhead} de cada hipervisor face ao sistema sem virtualização; e (iv)~um teste ao efeito do escalonador do Hyper-V.

O resto do artigo está organizado assim: a Secção~II apresenta o enquadramento e os trabalhos relacionados; a Secção~III, a arquitetura e os métodos; a Secção~IV, a configuração experimental; a Secção~V, os resultados; a Secção~VI, as limitações; e a Secção~VII, as conclusões.

% =============================================================== BACKGROUND
\section{Background e Trabalhos Relacionados}

\subsection{Virtualização x86 Assistida por Hardware}
Segundo Popek e Goldberg~\cite{b7}, uma arquitetura só é virtualizável se todas as instruções sensíveis forem privilegiadas, o que não acontecia no x86 original. Por isso, as primeiras soluções usavam tradução binária~\cite{b9} ou paravirtualização, como o Xen~\cite{b8}. Mais tarde, as extensões Intel VT-x e AMD-V criaram um modo de execução próprio para o hipervisor, e as tabelas de páginas de segundo nível (EPT/NPT) passaram a fazer em hardware a tradução dos endereços do convidado~\cite{b10}. Esta tradução tem um custo: quando há uma falha de TLB, é preciso percorrer as tabelas de páginas do convidado e as do hipervisor. Com páginas de 4~KiB, isto pode exigir até 24 acessos à memória; com páginas grandes (2~MiB), são muito menos~\cite{b10}. O outro custo importante são as saídas para o hipervisor (\textit{VM exits}), sobretudo quando o convidado acede a dispositivos.

Os dispositivos virtuais podem ser de dois tipos. Os \textit{emulados} imitam hardware real, como a placa de rede Intel \texttt{e1000} ou o controlador de disco LSI Logic, e cada acesso aos seus registos obriga a uma saída para o hipervisor. Os \textit{paravirtualizados} usam controladores que sabem que estão numa máquina virtual e trocam vários pedidos de uma vez através de memória partilhada, o que é mais eficiente~\cite{b8,b9}.

\subsection{Hipervisor Nativo: Microsoft Hyper-V}
No Hyper-V, o hipervisor (\texttt{hvix64}/\texttt{hvax64}) arranca antes do Windows, que passa a correr como \textit{partição raiz}. As máquinas virtuais correm em \textit{partições filhas} e usam dispositivos paravirtualizados: os controladores do convidado (VSC, como \texttt{hv\_storvsc} e \texttt{hv\_netvsc}) enviam os pedidos pelo barramento \textit{VMBus} para os fornecedores de serviço (VSP) da partição raiz, que os tratam em modo núcleo~\cite{b1}. O Hyper-V é um hipervisor de Tipo~1 tanto no Windows Server como no Windows~11. As edições diferem nas funcionalidades de centro de dados (como o \textit{Failover Clustering} e o Hyper-V Replica) e no escalonador usado por omissão: \textit{root} nas edições cliente, \textit{classic} no Windows Server~2016 e \textit{core} a partir do 2019~\cite{b11}. No \textit{root scheduler}, é o Windows que escalona os vCPU, como se fossem \textit{threads}. No \textit{core scheduler}, o hipervisor atribui núcleos físicos inteiros a cada máquina e nunca divide os \textit{threads} SMT de um núcleo entre máquinas diferentes, por razões de segurança~\cite{b11}.

\subsection{Hipervisor Baseado em Anfitrião: VMware Workstation}
O VMware Workstation instala um monitor (VMM) e um controlador no núcleo do Windows. No entanto, a emulação dos dispositivos e o acesso ao disco virtual (\texttt{.vmdk}) são feitos pelo processo \texttt{vmware-vmx}, em modo utilizador~\cite{b2,b9}. Por isso, a E/S do convidado chega ao NTFS através do Windows e pode ser servida pela \textit{cache} de ficheiros do anfitrião. Também a memória do convidado pertence ao processo \texttt{vmware-vmx} e é gerida pelo Windows. Quando o hipervisor do Windows está ativo (por exemplo, com o Hyper-V, a segurança baseada em virtualização ou o WSL2), o VMware não consegue usar o seu monitor nativo (modo \texttt{CPL0}) e passa a correr sobre a \textit{Windows Hypervisor Platform} (modo \texttt{ULM})~\cite{b12}. Nesse caso, deixa de ser um Tipo~2 ``puro''.

\subsection{Trabalhos Relacionados}
Adams e Agesen~\cite{b9} mostraram que as primeiras versões do VT-x eram, em algumas cargas, mais lentas do que a tradução binária, o que mudou com as EPT/NPT~\cite{b10}. Barham \textit{et~al.}~\cite{b8} mediram o custo da paravirtualização no Xen, e Felter \textit{et~al.}~\cite{b13} concluíram que, no KVM, o \textit{overhead} do processador e da memória é desprezável e que o maior custo está na E/S. Na Universidade da Beira Interior, Louro~\cite{b3} comparou uma infraestrutura de elevada disponibilidade com Hyper-V e outra com contentores Docker. Usou o IOzone (sistema de ficheiros), o RAMspeed (memória), o iPerf3 (rede) e o Geekbench (processador), em três ambientes por infraestrutura (nativo, virtualizado e aninhado), mas não indicou o número de repetições nem a dispersão dos resultados. Este trabalho avalia os mesmos subsistemas com ferramentas abertas e configuráveis, cujos resultados têm unidades físicas: o \texttt{sysbench}~\cite{b5} para o processador e a memória, o \texttt{fio}~\cite{b6} para o disco, com E/S direta e várias profundidades de fila, e o \texttt{iperf3} para a rede. Acrescenta ainda repetições e análise estatística.

% ================================================================ MÉTODOS
\section{Arquitetura e Métodos}

\subsection{Arquitetura dos Cenários}
A Fig.~\ref{fig_arq} mostra as duas pilhas de software. No Tipo~1, os pedidos de disco e de rede do convidado passam por controladores paravirtualizados até à partição raiz, onde são tratados em modo núcleo. No Tipo~2, os pedidos chegam a dispositivos emulados (LSI Logic e \texttt{e1000}, a configuração predefinida), e o processo \texttt{vmware-vmx} transforma-os em operações sobre ficheiros e sobre o adaptador virtual \texttt{VMnet8} do anfitrião.

\begin{figure}[t]
\centering
\resizebox{\columnwidth}{!}{%
\begin{tikzpicture}[
  font=\scriptsize,
  box/.style={draw, rounded corners=1.5pt, minimum height=5mm, align=center, inner sep=1.5pt},
  wide/.style={box, minimum width=37mm, text width=35mm},
  half/.style={box, minimum width=18mm, text width=16mm, minimum height=13mm, font=\tiny},
]
\node[wide, fill=gray!20] (hw1) at (0,0) {Hardware (i5-9300H, 16 GB, NVMe)};
\node[wide, fill=tipoum!30, above=1.2mm of hw1] (hyp) {Hipervisor Hyper-V};
\node[wide, fill=tipoum!10, minimum height=3.5mm, font=\tiny, above=1.2mm of hyp] (bus) {VMBus (memória partilhada)};
\node[half, fill=tipoum!10, above=1.2mm of bus.north west, anchor=south west] (root)
  {Partição raiz\\Windows 11\\StorVSP, VMSwitch\\VHDX $\rightarrow$ NTFS};
\node[half, fill=green!12, above=1.2mm of bus.north east, anchor=south east] (child)
  {Partição filha\\Ubuntu 26.04\\\texttt{hv\_storvsc}\\\texttt{hv\_netvsc}};
\node[above=0.8mm of bus.north, yshift=13mm, font=\footnotesize\bfseries, anchor=south] {Tipo 1 (nativo)};
\node[wide, fill=gray!20] (hw2) at (4.15,0) {Hardware (i5-9300H, 16 GB, NVMe)};
\node[wide, fill=tipodois!20, above=1.2mm of hw2] (os) {Windows 11: NTFS + \textit{cache}, VMnet8};
\node[wide, fill=tipodois!40, above=1.2mm of os, minimum height=3.5mm, font=\tiny] (vmx) {\texttt{vmware-vmx} (modo utilizador) + VMM (CPL0)};
\node[wide, fill=green!12, minimum height=13mm, above=1.2mm of vmx, font=\tiny] (g2) {Ubuntu 26.04 (convidado)\\LSI Logic (emulado)\\\texttt{e1000} (emulada)};
\node[above=0.8mm of g2, font=\footnotesize\bfseries] {Tipo 2 (\textit{host-based})};
\end{tikzpicture}}
\caption{Pilhas de software dos dois cenários e respetivos dispositivos virtuais.}
\label{fig_arq}
\end{figure}

\subsection{Formulação Estatística}
Cada métrica foi medida $N=10$ vezes. Para cada uma, calcula-se a média $\mu$, o desvio-padrão amostral
\begin{equation}
\sigma=\sqrt{\tfrac{1}{N-1}\textstyle\sum_{i=1}^{N}(x_i-\mu)^2},
\label{eq_sd}
\end{equation}
e o intervalo de confiança a 95\%, $\mu\pm t_{0{,}975;9}\,\sigma/\sqrt{N}$, com $t_{0{,}975;9}=2{,}262$. Para saber se a diferença entre dois cenários é significativa, usa-se o teste~$t$ de Welch, que não exige variâncias iguais, com o limite $p<0{,}05$.

\subsection{Quantificação do Overhead}
Tomando o Hyper-V como referência, a variação do Tipo~2 é
\begin{equation}
\Delta\% = \frac{\mu_{\mathrm{Tipo\,2}}-\mu_{\mathrm{Tipo\,1}}}{\mu_{\mathrm{Tipo\,1}}}\times 100 .
\label{eq_delta}
\end{equation}
Nas métricas de débito (eventos/s, MiB/s, IOPS, Mbit/s), um $\Delta\%$ negativo significa que o Tipo~2 é pior; nas de latência, o Tipo~2 é pior quando $\Delta\%$ é positivo.

\subsection{Modelo de Filas para o Armazenamento}
Segundo a lei de Little~\cite{b14}, quando há $Q$ pedidos em curso e a latência média é $\bar{L}$, o débito é
\begin{equation}
\mathrm{IOPS} = \frac{Q}{\bar{L}} .
\label{eq_little}
\end{equation}
Se o caminho de E/S consegue tratar vários pedidos em paralelo, aumentar $Q$ aumenta os IOPS sem aumentar muito a latência. Se só trata um pedido de cada vez, os IOPS deixam de subir e a latência cresce na mesma proporção que $Q$. Esta relação permite distinguir um caminho \textit{lento} de um caminho \textit{serializado}.

% ========================================================== CONFIGURAÇÃO
\section{Configuração Experimental}

\subsection{Sistema Anfitrião e Máquinas Virtuais}
Os dois cenários correram no mesmo portátil (Tabela~\ref{tab_host}), com Windows~11 Home~25H2 (\textit{build} 26200). A Tabela~\ref{tab_vm} resume as máquinas virtuais. Ambas têm os mesmos recursos e o mesmo sistema convidado, e confirmou-se que o \textit{kernel} e as versões das ferramentas eram iguais nos dois cenários.

\begin{table}[t]
\caption{Sistema Anfitrião}
\label{tab_host}
\centering
\footnotesize
\begin{tabular}{@{}ll@{}}
\toprule
\textbf{Componente} & \textbf{Especificação} \\
\midrule
Processador & Intel Core i5-9300H, 4 núcleos / 8 \textit{threads}, 2,4 GHz \\
Memória & 16 GB DDR4-2666 (2 $\times$ 8 GB) \\
Armazenamento & SSD NVMe Samsung PM981 (MZVLB1T0), 1 TB \\
Sistema operativo & Windows 11 Home 25H2, \textit{build} 26200 \\
Energia & Plano ``Elevado desempenho'', ligado à corrente \\
\bottomrule
\end{tabular}
\end{table}

\begin{table}[t]
\caption{Configuração das Máquinas Virtuais}
\label{tab_vm}
\centering
\scriptsize
\setlength{\tabcolsep}{2pt}
\begin{tabular}{@{}lll@{}}
\toprule
\textbf{Parâmetro} & \textbf{Hyper-V (Tipo 1)} & \textbf{VMware (Tipo 2)} \\
\midrule
Versão & config. VM 12.0, Geração 2 & Workstation 17.6.4, HW v21 \\
Convidado & \multicolumn{2}{l}{Ubuntu Server 26.04.1 LTS, \textit{kernel} 7.0.0-38} \\
vCPU & 2, sem SMT (1 \textit{socket} $\times$ 2) & 2, sem SMT (2 \textit{sockets} $\times$ 1) \\
Memória & 4096 MB (fixa) & 4096 MB \\
Disco (25 GB) & \texttt{.vhdx} dinâmico & \texttt{.vmdk} esparso \\
Controlador de disco & SCSI sintético (VMBus) & LSI Logic (emulado) \\
Escalonador E/S conv. & \texttt{none} & \texttt{mq-deadline} \\
Rede & \texttt{hv\_netvsc}, \textit{Default Switch} & \texttt{e1000}, NAT (VMnet8) \\
Sistema de ficheiros & \multicolumn{2}{l}{\texttt{ext4} sobre LVM} \\
Ferramentas & \multicolumn{2}{l}{\texttt{sysbench} 1.0.20, \texttt{fio} 3.41, \texttt{iperf3} 3.20} \\
Servidor de rede & \multicolumn{2}{l}{\texttt{iperf3} 3.22 no Windows anfitrião} \\
\bottomrule
\end{tabular}
\end{table}

\subsection{Instalação e Protocolo de Isolamento}
O Hyper-V é uma funcionalidade opcional do Windows (\texttt{Microsoft-Hyper-V-All}), mas não aparece na edição Home. Foi instalado com um \textit{script} que adiciona, com o DISM, os pacotes do Hyper-V que já vêm no Windows e depois ativa a funcionalidade, seguido de reinício~\cite{b15}. A máquina foi criada no Gestor Hyper-V como Geração~2, sem memória dinâmica. Antes dos testes, o \textit{checkpoint} automático foi fundido, para o disco não ficar como um VHDX de diferenciação. O VMware Workstation foi instalado com o instalador oficial, e a máquina foi criada com a configuração predefinida para Ubuntu de 64~bits. Nos dois casos, o Ubuntu Server~26.04 foi instalado com as opções predefinidas.

Por omissão, o Hyper-V mostra ao convidado os dois vCPU como dois \textit{threads} SMT do mesmo núcleo, enquanto o VMware os mostra como processadores independentes. Para a comparação depender só do hipervisor, desligou-se o SMT no Hyper-V (\texttt{Set-VMProcessor -HwThreadCountPerCore 1}). Assim, os dois convidados ficaram com dois vCPU independentes, sem SMT, num só nó NUMA, o que se confirmou com o \texttt{lscpu}. A única diferença que fica é a forma como os vCPU são agrupados: num \textit{socket} no Hyper-V e em dois no VMware. O Hyper-V não deixa definir \textit{sockets} sem criar nós NUMA virtuais, por isso esta diferença foi mantida. Não afeta o desempenho, porque o anfitrião só tem um processador físico.

Cada cenário precisou de um modo de arranque diferente do anfitrião, confirmado depois de cada reinício:
\begin{itemize}
\item \textbf{Hyper-V:} hipervisor ativo (\texttt{hypervisorlaunchtype auto}) e escalonador \textit{root} (tipo \texttt{0x4}, visto no registo de eventos do Windows).
\item \textbf{VMware:} hipervisor e Integridade de Memória desligados (\texttt{hypervisorlaunchtype off}), com \texttt{HypervisorPresent = False} e o ficheiro \texttt{vmware.log} a indicar \texttt{Monitor Mode: CPL0}.
\end{itemize}
Esta verificação é essencial: se o hipervisor do Windows estivesse ativo, o VMware correria em modo \texttt{ULM}, sobre o Hyper-V, e já não seria um Tipo~2.

Nos dois cenários: só uma máquina virtual ligada de cada vez; aplicações do anfitrião fechadas (incluindo a sincronização do OneDrive) e Windows Update em pausa; pastas das máquinas excluídas da análise em tempo real do Microsoft Defender; plano de energia ``Elevado desempenho'' e portátil sempre ligado à corrente. Antes de cada bateria de testes, esperou-se 10~minutos depois do arranque e suspenderam-se as atualizações automáticas do convidado (\texttt{apt}). Os testes foram lançados por SSH, dentro de uma sessão \texttt{tmux} no convidado.

\subsection{Cargas de Trabalho}
Os testes foram automatizados num só \textit{script}, igual nos dois convidados (Tabela~\ref{tab_cargas}). Cada bateria começa com uma corrida de aquecimento, que é descartada, seguida de $N=10$ repetições. Em cada repetição correm, por ordem, o processador, a memória e o disco, para que uma perturbação passageira do anfitrião não afete só um dos testes. Há pausas de 5~s entre testes, e a \textit{cache} de páginas do convidado é limpa antes de cada teste de disco (\texttt{drop\_caches}). Na rede, usou-se um servidor \texttt{iperf3} no Windows anfitrião, também com aquecimento e 10 repetições. Como a configuração é de nó único (\textit{single node}), o tráfego vai do convidado ao anfitrião apenas pelo comutador virtual de cada hipervisor, sem passar por uma placa de rede física. Assim, mede-se só o custo da virtualização da rede, que de outra forma ficaria escondido pelo limite de uma ligação física (cerca de 1~Gbit/s), igual nos dois cenários. Os comandos, os \textit{scripts} e os dados brutos estão no repositório do trabalho~\cite{b15}.

\textbf{Sistema de referência.} Para medir o \textit{overhead} face ao sistema sem virtualização, os testes foram repetidos no mesmo portátil, sem nenhuma máquina virtual, com as mesmas ferramentas e parâmetros. No processador e na memória, o \texttt{sysbench} correu num Ubuntu~26.04.1 nativo, arrancado a partir de uma \textit{pen} USB (\textit{live}, \textit{kernel} 7.0.0-30), com o perfil de energia ``desempenho''. Os dois \textit{threads} foram fixados em dois núcleos físicos diferentes (\texttt{taskset}), tal como os dois vCPU sem SMT das máquinas virtuais. O \textit{turbo} foi desligado (\texttt{intel\_pstate/no\_turbo}), o que limita o processador a 2,4~GHz. A razão é que, no Windows anfitrião, o processador corre sempre a cerca de 2,35~GHz, sem \textit{turbo}, em todos os planos de energia, e foi a essa frequência que correram as máquinas virtuais. Com o \textit{turbo}, o nativo chega a cerca de 4~GHz e é 67--71\% mais rápido no processador, o que mediria a frequência e não a virtualização. No disco e na rede, a referência foi o Windows anfitrião, com o hipervisor desligado (o mesmo modo do cenário VMware): o \texttt{fio}~3.41 com os seis perfis da Tabela~\ref{tab_cargas} (\texttt{windowsaio}, \texttt{direct=1}, ficheiro de 4~GiB) e o \texttt{iperf3}~3.22 em \textit{loopback} (127.0.0.1), que dá o limite máximo da pilha de rede do anfitrião. Em todos os casos, fizeram-se 10 repetições, com aquecimento.

\begin{table}[t]
\caption{Cargas de Trabalho (por repetição)}
\label{tab_cargas}
\centering
\scriptsize
\setlength{\tabcolsep}{2.5pt}
\begin{tabular}{@{}lll@{}}
\toprule
\textbf{Subsistema} & \textbf{Ferramenta e parâmetros} & \textbf{Duração} \\
\midrule
CPU & \texttt{sysbench cpu}, primos até 20\,000; 1 e 2 \textit{threads} & 10 s \\
\midrule
\multirow{4}{*}{Memória} & \texttt{sysbench memory}, escrita seq., bloco 1 KiB & 10 s \\
 & idem, escrita sequencial, bloco 1 MiB & 10 s \\
 & idem, leitura sequencial, bloco 1 MiB & 10 s \\
 & idem, leitura aleatória num \textit{buffer} de 256 MiB & 10 s \\
\midrule
\multirow{4}{*}{Disco} & \texttt{fio}, \texttt{libaio}, \texttt{direct=1}, ficheiro de 4 GiB: & \multirow{4}{*}{3+15 s} \\
 & \quad leitura e escrita aleatórias 4 KiB, $Q=1$ & \\
 & \quad mista 50/50 aleatória 4 KiB, $Q=1$ e $Q=32$ & \\
 & \quad leitura e escrita sequenciais 1 MiB, $Q=16$ & \\
\midrule
\multirow{3}{*}{Rede} & \texttt{iperf3} TCP: convidado$\to$anfitrião, inverso, 4 ligações & 10 s \\
 & \texttt{iperf3} UDP: 1 Mbit/s, 100 Mbit/s e 1 Gbit/s & 10 s \\
 & \texttt{ping}: 50 pacotes ICMP a 200 ms (RTT) & 10 s \\
\bottomrule
\end{tabular}
\end{table}

% ============================================================= RESULTADOS
\section{Resultados Experimentais}
A Tabela~\ref{tab_res} junta as métricas principais, e a Fig.~\ref{fig_rel} mostra o desempenho do Tipo~2 em relação ao Tipo~1. Todas as repetições foram mantidas, incluindo três valores fora do normal no VMware (uma escrita sequencial de 266~MiB/s, uma corrida de CPU com 1 \textit{thread} a 245~eventos/s e um RTT máximo de 10~ms). Estes valores aumentam o desvio-padrão, mas não mudam as conclusões.

\begin{table*}[t]
\caption{Resultados: Hyper-V (Tipo 1, escalonador \textit{root}) vs.\ VMware Workstation (Tipo 2), $\mu\pm\sigma$, $N=10$}
\label{tab_res}
\centering
\footnotesize
\begin{tabular}{@{}llrrrrc@{}}
\toprule
\textbf{Subsistema} & \textbf{Métrica} & \textbf{Hyper-V} & \textbf{VMware} & \textbf{$\Delta\%$} & \textbf{$p$ (Welch)} & \textbf{Melhor} \\
\midrule
\multirow{3}{*}{CPU}
 & Eventos/s, 1 \textit{thread} & 295,2 $\pm$ 2,9 & 287,9 $\pm$ 18,3 & $-2{,}5$ & 0,241 & n.s. \\
 & Eventos/s, 2 \textit{threads} & 552,8 $\pm$ 9,5 & 536,5 $\pm$ 22,3 & $-2{,}9$ & 0,055 & n.s. \\
 & Latência p95, 2 \textit{threads} (ms) & 4,44 $\pm$ 0,06 & 4,58 $\pm$ 0,25 & $+3{,}0$ & 0,132 & n.s. \\
\midrule
\multirow{4}{*}{Memória (MiB/s)}
 & Escrita seq., bloco 1 KiB & 4\,226 $\pm$ 267 & 4\,346 $\pm$ 163 & $+2{,}8$ & 0,244 & n.s. \\
 & Leitura seq., bloco 1 MiB & 29\,140 $\pm$ 859 & 28\,840 $\pm$ 2\,122 & $-1{,}0$ & 0,686 & n.s. \\
 & Escrita seq., bloco 1 MiB & 22\,499 $\pm$ 1\,181 & 22\,194 $\pm$ 2\,835 & $-1{,}4$ & 0,759 & n.s. \\
 & \textbf{Leitura aleatória, 256 MiB} & \textbf{608 $\pm$ 33} & \textbf{441 $\pm$ 37} & $\mathbf{-27{,}4}$ & $<$\textbf{0,001} & \textbf{Tipo 1} \\
\midrule
\multirow{7}{*}{Disco}
 & Leitura aleatória 4K, $Q=1$ (IOPS) & 3\,986 $\pm$ 183 & 6\,668 $\pm$ 1\,006 & $+67{,}3$ & $<$0,001 & Tipo 2 \\
 & \quad latência média / p99 ($\mu$s) & 249 / 356 & 151 / 182 & $-39{,}5$ / $-49{,}0$ & $<$0,001 / $<$0,001 & Tipo 2 \\
 & Escrita aleatória 4K, $Q=1$ (IOPS) & 5\,393 $\pm$ 295 & 6\,282 $\pm$ 770 & $+16{,}5$ & 0,006 & Tipo 2 \\
 & Mista 4K, $Q=1$ (IOPS totais) & 4\,656 $\pm$ 206 & 6\,183 $\pm$ 971 & $+32{,}8$ & $<$0,001 & Tipo 2 \\
 & \textbf{Mista 4K, $Q=32$ (IOPS totais)} & \textbf{75\,399 $\pm$ 2\,698} & \textbf{9\,161 $\pm$ 1\,152} & $\mathbf{-87{,}9}$ & $<$\textbf{0,001} & \textbf{Tipo 1} \\
 & Leitura sequencial 1M, $Q=16$ (MiB/s) & 1\,293 $\pm$ 31 & 972 $\pm$ 30 & $-24{,}9$ & $<$0,001 & Tipo 1 \\
 & Escrita sequencial 1M, $Q=16$ (MiB/s) & 1\,651 $\pm$ 232 & 798 $\pm$ 196 & $-51{,}7$ & $<$0,001 & Tipo 1 \\
\midrule
\multirow{7}{*}{Rede}
 & TCP convidado$\to$anfitrião (Mbit/s) & 3\,546 $\pm$ 241 & 2\,374 $\pm$ 172 & $-33{,}0$ & $<$0,001 & Tipo 1 \\
 & TCP anfitrião$\to$convidado (Mbit/s) & 6\,165 $\pm$ 539 & 2\,100 $\pm$ 73 & $-65{,}9$ & $<$0,001 & Tipo 1 \\
 & TCP 4 ligações (Mbit/s) & 4\,584 $\pm$ 767 & 2\,282 $\pm$ 143 & $-50{,}2$ & $<$0,001 & Tipo 1 \\
 & RTT médio / máximo (ms) & 0,46 / 0,78 & 0,51 / 2,33 & $+9{,}6$ / $+197$ & 0,287 / 0,129 & n.s. \\
 & UDP 1 Mbit/s: \textit{jitter} (ms) & 0,044 $\pm$ 0,027 & 0,008 $\pm$ 0,005 & $-82{,}7$ & 0,002 & Tipo 2 \\
 & UDP 100 Mbit/s: \textit{jitter} (ms) & 0,020 $\pm$ 0,010 & 0,126 $\pm$ 0,039 & $+547$ & $<$0,001 & Tipo 1 \\
 & UDP 1 Gbit/s: débito (Mbit/s) / perda (\%) & 977 / 2,3 & 299 / 0,0 & $-69{,}4$ / -- & $<$0,001 / 0,003 & Tipo 1 \\
\bottomrule
\multicolumn{7}{@{}l}{\scriptsize n.s.: diferença não significativa ($p\geq0{,}05$). Convenção de sinais de $\Delta\%$ conforme a Secção III-D.}
\end{tabular}
\end{table*}

\begin{figure}[t]
\centering
\begin{tikzpicture}
\begin{axis}[
  xbar, width=0.80\columnwidth, height=8.0cm,
  xmin=0, xmax=215, bar width=5pt,
  xlabel={VMware relativo ao Hyper-V (\%)},
  symbolic y coords={Rede UDP 1G,Rede 4 lig.,Rede anf.$\to$conv.,Rede conv.$\to$anf.,Disco seq. escrita,Disco seq. leitura,Disco mista Q32,Disco mista Q1,Disco leit. aleat. Q1,Mem. aleat. 256M,Mem. escrita 1M,Mem. leitura 1M,Mem. escrita 1K,CPU 2 threads,CPU 1 thread},
  ytick=data, y tick label style={font=\tiny},
  x tick label style={font=\scriptsize}, xlabel style={font=\scriptsize},
  nodes near coords, nodes near coords align={horizontal},
  every node near coord/.append style={font=\tiny, xshift=13pt, fill=white, inner sep=0.5pt},
  /pgf/number format/.cd, fixed, precision=0,
  enlarge y limits=0.05,
  extra x ticks={100}, extra x tick style={grid=major, grid style={dashed, black!60}, tick label style={font=\scriptsize\bfseries}},
]
\addplot[fill=tipodois!70, draw=tipodois!90!black, error bars/.cd, x dir=both, x explicit]
  coordinates {
    (97.5,CPU 1 thread) +- (6.2,0)
    (97.1,CPU 2 threads) +- (4.0,0)
    (102.8,Mem. escrita 1K) +- (3.9,0)
    (99.0,Mem. leitura 1M) +- (7.3,0)
    (98.6,Mem. escrita 1M) +- (12.6,0)
    (72.6,Mem. aleat. 256M) +- (6.1,0)
    (167.3,Disco leit. aleat. Q1) +- (25.2,0)
    (132.8,Disco mista Q1) +- (20.9,0)
    (12.2,Disco mista Q32) +- (1.5,0)
    (75.1,Disco seq. leitura) +- (2.3,0)
    (48.3,Disco seq. escrita) +- (11.9,0)
    (67.0,Rede conv.$\to$anf.) +- (4.9,0)
    (34.1,Rede anf.$\to$conv.) +- (1.2,0)
    (49.8,Rede 4 lig.) +- (3.1,0)
    (30.6,Rede UDP 1G) +- (1.0,0)
  };
\end{axis}
\end{tikzpicture}
\caption{Débito do Tipo~2 relativo ao Tipo~1 (Hyper-V $=100\%$, tracejado). Barras de erro: $\sigma_{\mathrm{VMware}}/\mu_{\mathrm{Hyper\text{-}V}}$ (não incluem a dispersão do Hyper-V, indicada na Tabela~\ref{tab_res}).}
\label{fig_rel}
\end{figure}

\subsection{Processador}
Com um \textit{thread}, o Hyper-V obteve 295,2~eventos/s e o VMware 287,9 ($-2{,}5\%$, $p=0{,}241$). Com dois \textit{threads}, obtiveram 552,8 e 536,5 ($-2{,}9\%$, $p=0{,}055$). Nenhuma das diferenças é significativa, e os dois escalam quase da mesma forma de 1 para 2 \textit{threads} (1,87 no Hyper-V e 1,86 no VMware). Isto está de acordo com a literatura: com virtualização assistida por hardware, o código do convidado corre diretamente no processador nas duas arquiteturas~\cite{b9,b13}, e fica a menos de 5\% do nativo com um \textit{thread} (Tabela~\ref{tab_nat}). A diferença está na \textit{estabilidade}. O coeficiente de variação foi de 1,0--1,7\% no Hyper-V e de 4,2--6,4\% no VMware. Uma explicação provável é que, no Tipo~2, os vCPU são \textit{threads} de um processo normal e competem com as outras tarefas do Windows.

\textbf{Efeito do escalonador.} O teste de processador foi repetido com o \textit{core scheduler} (\texttt{hypervisorschedulertype core}), o escalonador por omissão do Windows Server, com a mesma topologia (Fig.~\ref{fig_cpu}). Com dois \textit{threads}, o resultado foi igual (555,6 contra 552,8~eventos/s, $+0{,}5\%$, $p=0{,}77$). Com um \textit{thread}, foi 6,1\% mais baixo (277,2 contra 295,2, $p<0{,}001$) e mais irregular (percentil~95 de 4,56 contra 3,59~ms). Ou seja, numa máquina com vCPU independentes, o escalonador não muda o débito total. A pequena perda com um \textit{thread} pode dever-se ao trabalho extra que o \textit{core scheduler} tem para manter cada núcleo físico isolado~\cite{b11}. Assim, o escalonador do Windows~11 não dá vantagem ao Hyper-V em relação ao do Windows Server.

\begin{figure}[t]
\centering
\begin{tikzpicture}
\begin{axis}[
  ybar, width=0.95\columnwidth, height=4.3cm,
  bar width=10pt, ymin=0, ymax=680,
  ylabel={Eventos por segundo},
  symbolic x coords={1 thread,2 threads}, xtick=data,
  enlarge x limits=0.4,
  legend style={at={(0.5,1.03)}, anchor=south, font=\tiny, legend columns=3},
  tick label style={font=\scriptsize}, label style={font=\scriptsize},
  nodes near coords, every node near coord/.append style={font=\tiny, yshift=4pt},
  /pgf/number format/.cd, fixed, precision=0,
]
\addplot[fill=tipoum!80, draw=tipoum!90!black, error bars/.cd, y dir=both, y explicit]
  coordinates {(1 thread,295.2) +- (0,2.9) (2 threads,552.8) +- (0,9.5)};
\addplot[fill=tipocore, draw=tipoum!70!black, error bars/.cd, y dir=both, y explicit]
  coordinates {(1 thread,277.2) +- (0,10.5) (2 threads,555.6) +- (0,28.7)};
\addplot[fill=tipodois!70, draw=tipodois!90!black, error bars/.cd, y dir=both, y explicit]
  coordinates {(1 thread,287.9) +- (0,18.3) (2 threads,536.5) +- (0,22.3)};
\legend{Hyper-V \textit{root}, Hyper-V \textit{core}, VMware}
\end{axis}
\end{tikzpicture}
\caption{Processador ($\mu\pm\sigma$): paridade entre os três cenários com dois \textit{threads}.}
\label{fig_cpu}
\end{figure}

\subsection{Memória}\label{sec_mem}
Nos três testes sequenciais, as diferenças vão de $-1{,}4\%$ a $+2{,}8\%$ e nenhuma é significativa: a largura de banda da RAM (blocos de 1~MiB, maiores do que as \textit{caches}) e o acesso dentro da \textit{cache} L1 (blocos de 1~KiB) são iguais nos dois. Já o acesso \textit{aleatório} num \textit{buffer} de 256~MiB é 27,4\% mais lento no VMware (441 contra 608~MiB/s, $p<0{,}001$). É o único teste de memória que provoca falhas de TLB constantes e, por isso, obriga a percorrer as tabelas de páginas do convidado e do hipervisor. Este custo depende do tamanho das páginas que o hipervisor usa para a memória do convidado~\cite{b10}. A diferença é compatível com o Hyper-V usar páginas grandes, enquanto no VMware a memória do convidado pertence ao processo \texttt{vmware-vmx} e o Windows a atribui, por omissão, em páginas de 4~KiB. Em resumo, no acesso sequencial a memória não tem \textit{overhead} relativo (embora os dois fiquem 12--20\% abaixo do nativo), mas no Tipo~2 há um custo claro em cargas com acessos dispersos, como bases de dados.

\subsection{Armazenamento}
No disco, o resultado depende do número de pedidos em fila.

\textbf{Um pedido de cada vez ($Q=1$): vantagem do Tipo~2.} Com um pedido de cada vez, o VMware teve mais 67\% de IOPS em leitura aleatória, com latência média de 151~$\mu$s contra 249~$\mu$s e um percentil~99 de 182~$\mu$s contra 356~$\mu$s. A explicação mais provável é a \textit{cache} de ficheiros do Windows. O ficheiro de teste de 4~GiB cabe nos 16~GB do anfitrião, e o \texttt{vmware-vmx} lê o \texttt{.vmdk} através do sistema de ficheiros. Assim, as leituras ``diretas'' do convidado acabam por ser servidas pela RAM do anfitrião: a opção \texttt{direct=1} evita a \textit{cache} do convidado, mas não a do anfitrião. A escrita beneficia do mesmo efeito ($+16{,}5\%$ IOPS), porque o convidado recebe a confirmação antes de os dados chegarem ao SSD. Isto melhora a latência, mas põe os dados em risco se houver uma falha de energia. No Hyper-V, cada pedido vai pelo VMBus até ao NVMe, sem esta \textit{cache}.

\textbf{Vários pedidos em fila: vantagem do Tipo~1.} Com 32 pedidos em fila, o Hyper-V passou de 4\,656 para 75\,399~IOPS (16 vezes mais), com uma latência média (entre leituras e escritas) de apenas 424~$\mu$s. O VMware passou de 6\,183 para 9\,161 (1,5 vezes), com uma latência de 3\,541~$\mu$s, ou seja, $-88\%$ ($p<0{,}001$). A lei de Little~\eqref{eq_little} confirma os dois valores: $32/424\,\mu\mathrm{s}\approx75\,500$ e $32/3\,541\,\mu\mathrm{s}\approx9\,036$~IOPS, a menos de 1,5\% do medido. No VMware, a latência cresce na proporção da fila e o débito fica preso em cerca de 9\,000~IOPS. O caminho de E/S está, portanto, \textit{serializado}, o que é de esperar num controlador emulado (LSI Logic) tratado por um processo em modo utilizador. O controlador do Hyper-V, pelo contrário, envia os pedidos em paralelo para o NVMe. O mesmo limite vê-se no débito sequencial com blocos de 1~MiB: $-25\%$ em leitura e $-52\%$ em escrita no Tipo~2 (Fig.~\ref{fig_disco}).

\begin{figure}[t]
\centering
\begin{tikzpicture}
\begin{axis}[
  ybar, width=0.95\columnwidth, height=4.4cm,
  bar width=13pt, ymode=log, ymin=1000, ymax=200000,
  log origin=infty,
  ylabel={IOPS totais (escala log.)},
  symbolic x coords={$Q=1$,$Q=32$}, xtick=data,
  enlarge x limits=0.45,
  legend style={at={(0.03,0.97)}, anchor=north west, font=\scriptsize},
  tick label style={font=\scriptsize}, label style={font=\scriptsize},
  nodes near coords, every node near coord/.append style={font=\tiny, yshift=2pt},
  /pgf/number format/.cd, fixed, precision=0, 1000 sep={\,},
  point meta=rawy,
]
\addplot[fill=tipoum!80, draw=tipoum!90!black] coordinates {($Q=1$,4656) ($Q=32$,75399)};
\addplot[fill=tipodois!70, draw=tipodois!90!black] coordinates {($Q=1$,6183) ($Q=32$,9161)};
\legend{Hyper-V (Tipo 1), VMware (Tipo 2)}
\end{axis}
\end{tikzpicture}
\caption{E/S aleatória mista 4~KiB: o Tipo~2 é mais rápido com $Q=1$ (\textit{cache} do anfitrião), mas não escala com a profundidade de fila.}
\label{fig_disco}
\end{figure}

\subsection{Rede}
Em TCP, o Hyper-V foi 33\% mais rápido do convidado para o anfitrião, 66\% mais rápido no sentido inverso e 50\% mais rápido com quatro ligações em paralelo, sempre com $p<0{,}001$ (Fig.~\ref{fig_rede}). No sentido inverso, o \texttt{iperf3} do anfitrião chegou a 87\% de CPU, por isso o valor do Hyper-V pode estar limitado pelo próprio gerador de tráfego. A latência ida-e-volta média é igual (0,46 contra 0,51~ms, $p=0{,}29$), mas o VMware tem picos até 2,3~ms e uma variação quatro vezes maior. A causa é o tipo de placa de rede. A \texttt{e1000} do VMware é emulada: cada acesso aos seus registos obriga a sair para o monitor e a processar no \texttt{vmware-vmx}. A \texttt{hv\_netvsc} do Hyper-V troca pacotes em lote pelo VMBus. Ao enviar a partir do convidado, o \texttt{iperf3} do anfitrião gastou 31,6\% de CPU no VMware e 9,4\% no Hyper-V, com um débito mais baixo: cerca de 75 contra 376~Mbit/s por ponto percentual de CPU. No VMware, cada bit recebido custa mais CPU ao anfitrião. Com a placa paravirtualizada \texttt{vmxnet3}, a diferença seria provavelmente menor, mas não foi testada, para manter a configuração predefinida (Secção~VI). Os testes UDP mostram o mesmo. Com pouco tráfego (1~Mbit/s), o \textit{jitter} é muito baixo nos dois (0,044 contra 0,008~ms). Com 100~Mbit/s, o \textit{jitter} do VMware é seis vezes maior (0,126 contra 0,020~ms, $p<0{,}001$). A 1~Gbit/s, o Hyper-V entrega 977~Mbit/s com 2,3\% de perda, e o VMware não passa dos 299~Mbit/s. O VMware não perde pacotes só porque não consegue enviar mais: o limite está na própria placa emulada, que envia cerca de um terço dos pacotes UDP do Tipo~1.

Tal como Louro~\cite{b3}, usou-se o \texttt{iperf3} em TCP, para o débito, e em UDP, para o \textit{jitter}. No trabalho de Louro, porém, o tráfego passava pela rede física entre os dois nós de um \textit{cluster}, com uma janela TCP fixa de 20~KB. As máquinas sobre Hyper-V chegaram a 207--222~Mbit/s e a um \textit{jitter} de 0,39--1,07~ms~\cite[Figs.~4.3--4.4]{b3}. Como os dois trabalhos medem caminhos diferentes, os valores absolutos não são comparáveis.

\begin{figure}[t]
\centering
\begin{tikzpicture}
\begin{axis}[
  ybar, width=0.95\columnwidth, height=4.3cm,
  bar width=11pt, ymin=0, ymax=8800,
  ylabel={Débito TCP (Mbit/s)},
  symbolic x coords={Conv.$\to$anf.,Anf.$\to$conv.,4 ligações}, xtick=data,
  enlarge x limits=0.25,
  legend style={at={(0.98,0.98)}, anchor=north east, font=\scriptsize},
  tick label style={font=\scriptsize}, label style={font=\scriptsize},
  nodes near coords, every node near coord/.append style={font=\tiny, yshift=4pt},
  /pgf/number format/.cd, fixed, precision=0, 1000 sep={\,},
]
\addplot[fill=tipoum!80, draw=tipoum!90!black, error bars/.cd, y dir=both, y explicit]
  coordinates {(Conv.$\to$anf.,3546) +- (0,241) (Anf.$\to$conv.,6165) +- (0,539) (4 ligações,4584) +- (0,767)};
\addplot[fill=tipodois!70, draw=tipodois!90!black, error bars/.cd, y dir=both, y explicit]
  coordinates {(Conv.$\to$anf.,2374) +- (0,172) (Anf.$\to$conv.,2100) +- (0,73) (4 ligações,2282) +- (0,143)};
\legend{Hyper-V (\texttt{hv\_netvsc}), VMware (\texttt{e1000})}
\end{axis}
\end{tikzpicture}
\caption{Débito de rede convidado$\leftrightarrow$anfitrião ($\mu\pm\sigma$).}
\label{fig_rede}
\end{figure}

\subsection{Overhead Face ao Sistema Nativo}\label{sec_nat}
A Tabela~\ref{tab_nat} compara cada hipervisor com o sistema de referência, usando a Eq.~\eqref{eq_delta} com o nativo no lugar do Hyper-V. No processador, o \textit{overhead} é pequeno nos dois: $-1{,}9\%$ (Hyper-V) e $-4{,}4\%$ (VMware) com um \textit{thread}, e $-8{,}2\%$ e $-10{,}9\%$ com dois. O nativo escala perfeitamente de 1 para 2 \textit{threads} (fator 2,00), contra 1,87 e 1,86 nas máquinas virtuais, e tem uma latência p95 de 3,30~ms, contra 4,44 e 4,58~ms. Assim, o custo de virtualizar o processador está sobretudo em correr os dois vCPU ao mesmo tempo, em competição com os \textit{threads} do anfitrião, e não em cada instrução. Cerca de 2~pontos percentuais vêm da diferença de frequência (2,35 contra 2,40~GHz).

Na memória sequencial com blocos de 1~MiB, os dois hipervisores perdem 12--20\% face ao nativo, sem diferença significativa entre eles. No acesso aleatório, o \textit{overhead} é de 26,8\% no Hyper-V e de 46,9\% no VMware. É o maior custo no processador e na memória, o que confirma que a tradução de endereços em dois níveis é o principal custo da virtualização da memória, e que é maior no Tipo~2 (Secção~\ref{sec_mem}). O teste com blocos de 1~KiB foi \textit{mais rápido} nas máquinas virtuais (+16 a +20\%). Neste teste, cada bloco conta como um evento, e o resultado depende do custo fixo que o \texttt{sysbench} tem por evento, que muda de ambiente para ambiente. Por isso, não é visto como um ganho da virtualização.

No disco, a referência é o Windows anfitrião, que usa outra pilha de E/S (NTFS e \texttt{windowsaio}, em vez de \texttt{ext4} e \texttt{libaio}), por isso a comparação serve apenas de indicação. Com $Q=1$, os dois hipervisores perdem 14--65\% de IOPS face ao anfitrião, e a latência média de leitura do Hyper-V é quase o dobro da nativa (249 contra 127~$\mu$s). É o custo, em cada pedido, de passar pelo VMBus e pela partição raiz. O VMware perde menos em leitura ($-14\%$) porque parte das leituras vem da \textit{cache} do anfitrião. Com $Q=32$ e em leitura sequencial, o Hyper-V fica \textit{acima} da referência ($+38\%$ e $+6\%$). Não é a virtualização que acelera o disco: o \texttt{libaio} do Linux aguenta mais pedidos em simultâneo do que o \texttt{windowsaio} do \texttt{fio} no Windows. O VMware, pelo contrário, fica a 17\% do anfitrião com $Q=32$ e a 42\% em escrita sequencial. Na rede, o \textit{loopback} do anfitrião chegou a 18\,084~Mbit/s em TCP, 17\,614 no sentido inverso e 45\,070 com quatro ligações. As máquinas virtuais chegam a 10--35\% deste limite no Hyper-V e a 5--13\% no VMware. Como o \textit{loopback} não passa por nenhum dispositivo de rede, estes valores mostram apenas a margem que cada comutador virtual deixa por aproveitar.

\begin{table}[t]
\caption{Overhead Face ao Sistema Nativo ($\mu$, $N=10$)}
\label{tab_nat}
\centering
\scriptsize
\setlength{\tabcolsep}{3pt}
\begin{tabular}{@{}lrrr@{}}
\toprule
\textbf{Métrica} & \textbf{Nativo} & \textbf{Hyper-V} & \textbf{VMware} \\
\midrule
\multicolumn{4}{@{}l}{\textit{Linux nativo, 2,4 GHz sem turbo (sysbench)}} \\
CPU, 1 \textit{thread} (eventos/s) & 301,0 & $-1{,}9\%$ & $-4{,}4\%$ \\
CPU, 2 \textit{threads} (eventos/s) & 602,1 & $-8{,}2\%$ & $-10{,}9\%$ \\
CPU, latência p95, 2 \textit{threads} (ms) & 3,30 & $+34{,}6\%$ & $+38{,}7\%$ \\
Mem.\ escrita seq.\ 1 KiB (MiB/s) & 3\,633 & $+16{,}3\%$ & $+19{,}6\%$ \\
Mem.\ leitura seq.\ 1 MiB (MiB/s) & 33\,199 & $-12{,}2\%$ & $-13{,}1\%$ \\
Mem.\ escrita seq.\ 1 MiB (MiB/s) & 27\,882 & $-19{,}3\%$ & $-20{,}4\%$ \\
\textbf{Mem.\ leitura aleat.\ 256 MiB (MiB/s)} & \textbf{831} & $\mathbf{-26{,}8\%}$ & $\mathbf{-46{,}9\%}$ \\
\midrule
\multicolumn{4}{@{}l}{\textit{Windows anfitrião, hipervisor desligado (fio, windowsaio)}} \\
Leitura aleat.\ 4K, $Q=1$ (IOPS) & 7\,758 & $-48{,}6\%$ & $-14{,}1\%$ \\
\quad latência média ($\mu$s) & 127 & $+95{,}5\%$ & $+18{,}3\%$ \\
Escrita aleat.\ 4K, $Q=1$ (IOPS) & 15\,382 & $-64{,}9\%$ & $-59{,}2\%$ \\
Mista 4K, $Q=1$ (IOPS totais) & 10\,259 & $-54{,}6\%$ & $-39{,}7\%$ \\
Mista 4K, $Q=32$ (IOPS totais) & 54\,494 & $+38{,}4\%$ & $-83{,}2\%$ \\
Leitura seq.\ 1M, $Q=16$ (MiB/s) & 1\,218 & $+6{,}2\%$ & $-20{,}2\%$ \\
Escrita seq.\ 1M, $Q=16$ (MiB/s) & 1\,910 & $-13{,}5\%$ & $-58{,}2\%$ \\
\midrule
\multicolumn{4}{@{}l}{\textit{Windows anfitrião, iperf3 em loopback (indicativo)}} \\
TCP envio/inverso (Mbit/s) & 18\,084/17\,614 & $-80$/$-65\%$ & $-87$/$-88\%$ \\
TCP 4 ligações (Mbit/s) & 45\,070 & $-89{,}8\%$ & $-94{,}9\%$ \\
\bottomrule
\multicolumn{4}{@{}p{0.95\columnwidth}@{}}{\scriptsize Todas as diferenças face ao nativo são significativas ($p<0{,}05$, Welch). Desvios-padrão e IC95 no ficheiro de resultados~\cite{b15}.}
\end{tabular}
\end{table}

% ============================================================ LIMITAÇÕES
\section{Limitações e Ameaças à Validade}
\begin{itemize}
\item \textbf{Referência:} não há um único sistema nativo equivalente às duas máquinas virtuais. O processador e a memória foram comparados com um Linux nativo arrancado por USB (\textit{kernel} 7.0.0-30, contra 7.0.0-38 nas máquinas virtuais), a 2,4~GHz fixos. O disco e a rede foram comparados com o Windows anfitrião (NTFS e \texttt{windowsaio}; e \textit{loopback}, porque a ligação convidado$\leftrightarrow$anfitrião não existe sem virtualização). Por isso, no disco e na rede, o \textit{overhead} face ao nativo inclui as diferenças entre sistemas operativos, e não só o custo do hipervisor.
\item \textbf{Dispositivos predefinidos:} o VMware usou dispositivos emulados (LSI Logic, \texttt{e1000}) e o Hyper-V paravirtualizados. É a configuração por omissão de cada produto, mas aumenta a diferença em relação a um VMware com \texttt{vmxnet3} e PVSCSI/NVMe. Os escalonadores de E/S do convidado também foram diferentes (\texttt{none} e \texttt{mq-deadline}).
\item \textbf{\textit{Cache} do anfitrião no Tipo~2:} o ficheiro de teste de 4~GiB cabe na memória do anfitrião, e o \texttt{direct=1} só evita a \textit{cache} do convidado. Assim, com $Q=1$, o VMware é medido em parte a partir da RAM do anfitrião e o Hyper-V a partir do NVMe. Os dois caminhos não são equivalentes, e a vantagem do Tipo~2 com $Q=1$ não quer dizer que o seu disco virtual seja mais rápido.
\item \textbf{Modo do anfitrião:} no cenário Hyper-V, o Windows correu como partição raiz, com a segurança baseada em virtualização ativa; no cenário VMware, correu sem hipervisor. Esta diferença é inevitável quando as duas arquiteturas estão no mesmo equipamento.
\item \textbf{Causa da diferença na memória:} a explicação para a leitura aleatória (tamanho das páginas) baseia-se no padrão dos resultados e na literatura, mas não foi medida diretamente.
\item \textbf{Equipamento:} um portátil com Windows~11 Home (edição em que o Hyper-V não é suportado oficialmente), em nó único e com 2~vCPU e 4~GB por máquina. Os resultados podem não se aplicar a servidores nem a infraestruturas com vários nós, onde a rede física também conta.
\item \textbf{\textit{Benchmarks} sintéticos:} medem subsistemas isolados e não cargas reais, que os combinam.
\end{itemize}

% ============================================================= CONCLUSÕES
\section{Conclusões}
Este trabalho comparou um hipervisor nativo (Hyper-V) e um hipervisor \textit{host-based} (VMware Workstation~17.6) no mesmo equipamento, com máquinas virtuais Ubuntu Server~26.04 com os mesmos recursos e vCPU independentes. Nas condições testadas:
\begin{enumerate}
\item \textbf{O processador e a memória sequencial} não têm diferenças significativas entre os dois hipervisores, graças à virtualização assistida por hardware. O Hyper-V foi, no entanto, mais estável (coeficiente de variação abaixo de 2\%).
\item \textbf{O acesso aleatório à memória} é 27\% mais lento no Tipo~2, o que é compatível com um maior custo da tradução de endereços em dois níveis.
\item \textbf{No disco}, o Tipo~2 é mais rápido com pedidos isolados, porque a \textit{cache} do anfitrião serve as leituras e confirma as escritas antes de chegarem ao SSD. Com vários pedidos em simultâneo, porém, fica limitado a cerca de 9\,000~IOPS, como a lei de Little confirma, e perde 88\% de IOPS com $Q=32$ e 25--52\% de débito sequencial.
\item \textbf{Na rede}, o Tipo~1 é 33--66\% mais rápido em TCP, consegue três vezes mais débito em UDP e tem seis vezes menos \textit{jitter} com carga moderada, graças aos dispositivos paravirtualizados.
\item \textbf{Face ao sistema nativo}, à mesma frequência, o processador perde 2--11\% nos dois hipervisores, sobretudo quando os dois vCPU trabalham ao mesmo tempo, e a memória sequencial perde 12--20\%. O maior custo está no acesso aleatório à memória: 27\% no Tipo~1 e 47\% no Tipo~2.
\item \textbf{O escalonador do Hyper-V} (\textit{root} ou \textit{core}) não muda o débito de uma máquina com vCPU independentes.
\end{enumerate}
Para cargas típicas de centro de dados, com muita E/S em simultâneo, rede intensiva e acessos dispersos à memória, os resultados favorecem claramente o hipervisor nativo. O Tipo~2 continua a ser adequado para computadores de trabalho, com cargas interativas e pouco paralelas. Num anfitrião Windows, é também essencial confirmar que o hipervisor de Tipo~2 corre no seu modo nativo (\texttt{CPL0}) e não sobre o Hyper-V. Como trabalho futuro, propõe-se repetir os testes com dispositivos paravirtualizados no VMware e avaliar cargas reais, como bases de dados e servidores web.

% ============================================================ REFERÊNCIAS
\begin{thebibliography}{00}
\bibitem{b1} Microsoft, ``Hyper-V overview,'' [Online]. Available: \url{https://learn.microsoft.com/pt-pt/windows-server/virtualization/hyper-v/overview}. [Accessed: Oct.~18, 2026].
\bibitem{b2} VMware, Inc., ``VMware Workstation Pro documentation,'' [Online]. Available: \url{https://www.vmware.com}. [Accessed: Oct.~18, 2026].
\bibitem{b3} A. Louro, ``Comparação do desempenho de infraestruturas virtualizadas de elevada disponibilidade usando hypervisors nativos ou containers,'' M.Sc. dissertation, Dept. Eng. Informática, Univ. Beira Interior, Covilhã, Portugal, 2018. [Online]. Available: \url{https://ubibliorum.ubi.pt/handle/10400.6/9837}.
\bibitem{b4} Universidade da Beira Interior, ``Conteúdos das aulas teóricas da unidade curricular Virtualização e Centros de Dados,'' [Online]. Available: \url{https://moodle.ubi.pt/}. [Accessed: Oct.~18, 2026].
\bibitem{b5} A. Kopytov, ``sysbench: Scriptable database and system performance benchmark,'' GitHub, 2020. [Online]. Available: \url{https://github.com/akopytov/sysbench}.
\bibitem{b6} J. Axboe, ``Flexible I/O Tester (fio),'' GitHub. [Online]. Available: \url{https://github.com/axboe/fio}. [Accessed: Oct.~18, 2026].
\bibitem{b7} G. J. Popek and R. P. Goldberg, ``Formal requirements for virtualizable third generation architectures,'' \textit{Commun. ACM}, vol.~17, no.~7, pp.~412--421, 1974.
\bibitem{b8} P. Barham \textit{et al.}, ``Xen and the art of virtualization,'' in \textit{Proc. 19th ACM Symp. Operating Systems Principles (SOSP)}, 2003, pp.~164--177.
\bibitem{b9} K. Adams and O. Agesen, ``A comparison of software and hardware techniques for x86 virtualization,'' in \textit{Proc. 12th Int. Conf. ASPLOS}, 2006, pp.~2--13.
\bibitem{b10} R. Bhargava, B. Serebrin, F. Spadini, and S. Manne, ``Accelerating two-dimensional page walks for virtualized systems,'' in \textit{Proc. 13th Int. Conf. ASPLOS}, 2008, pp.~26--35.
\bibitem{b11} Microsoft, ``Manage Hyper-V hypervisor scheduler types,'' [Online]. Available: \url{https://learn.microsoft.com/windows-server/virtualization/hyper-v/manage/manage-hyper-v-scheduler-types}. [Accessed: Oct.~22, 2026].
\bibitem{b12} Microsoft, ``Windows Hypervisor Platform,'' [Online]. Available: \url{https://learn.microsoft.com/virtualization/api/}. [Accessed: Oct.~22, 2026].
\bibitem{b13} W. Felter, A. Ferreira, R. Rajamony, and J. Rubio, ``An updated performance comparison of virtual machines and Linux containers,'' in \textit{Proc. IEEE ISPASS}, 2015, pp.~171--172.
\bibitem{b14} J. D. C. Little, ``A proof for the queuing formula: $L=\lambda W$,'' \textit{Operations Research}, vol.~9, no.~3, pp.~383--387, 1961.
\bibitem{b15} H. Laia, A. Santos e D. Basílio, ``Hyper-V vs. VMware: comandos, \textit{scripts} e dados experimentais,'' GitHub, 2026. [Online]. Available: \url{https://github.com/henriquelaia/vcd-hyperv-vs-vmware}.
\end{thebibliography}

\end{document}